\setcounter{section}{6}

\Needspace{5\baselineskip}
\hypertarget{klux6563ux5ea6ux4e0ejsdux6563ux5ea6}{%
\section{KL散度与JSD散度}\label{klux6563ux5ea6ux4e0ejsdux6563ux5ea6}}

在大模型中，KL散度衡量输出概率分布的差异，并非按输出前的隐藏层维度计算，而是在大小为 \(V\) 的词表上计算（因为本质上大模型的输出是离散的tokens）。

\[
D(P\Vert Q)=\sum_{i=1}^{V}p_i\log\frac{p_i}{q_i}
\]

注意：计算困惑度时是对序列中的所有token的概率负对数取平均，而计算KL散度时是在输出一个token时，对词表中的所有词的KL散度取平均，二者取平均的方式不同。

\Needspace{5\baselineskip}
实际研究中，为了让数值更稳定，往往用KL散度的变体------JSD散度，即取 \(P\) 和 \(Q\) 的平均 \(M\)：

\[
M=\frac{1}{2}(P+Q)
\]

\Needspace{5\baselineskip}
JSD散度定义为：

\[
\mathrm{JSD}(P,Q)=\frac{1}{2}\left(D_{\mathrm{KL}}(P\Vert M)+D_{\mathrm{KL}}(Q\Vert M)\right)
\]

其功能和KL类似但不会出现 \(\log\) 内分母为零，例如研究模型中间隐藏层和输出层之间神经元激活值的JSD散度（如果二者维度不同，则投影后再用JSD散度），就可以判断二者分布之间的相似程度大小。

\Needspace{5\baselineskip}
对于LLM，自回归输出 \(n\) 个token的KL散度：

\[
D_{\mathrm{KL}}(P(Y)\Vert Q(Y))
=
\sum_Y P(Y)\log\frac{P(Y)}{Q(Y)}.
\]

\Needspace{5\baselineskip}
由于直接对所有可能的序列求和是不可行的（有 \(V^n\) 种可能），需利用自回归特性将其分解：

\[
\begin{aligned}
D_{\mathrm{KL}}(P(Y)\Vert Q(Y))
&=\mathbb{E}_{Y\sim P}
\left[
\sum_{t=1}^{n}
\left(
\log P(y_t\mid y_{<t})-\log Q(y_t\mid y_{<t})
\right)
\right] \\
&=\sum_{t=1}^{n}
\mathbb{E}_{Y\sim P}
\left[
\log\frac{P(y_t\mid y_{<t})}{Q(y_t\mid y_{<t})}
\right].
\end{aligned}
\]

\Needspace{5\baselineskip}
这里的处理难点在于 \(Y\sim P\)，需要考虑整个序列的概率分布。注意到右边的条件概率的条件为 \(y_{<t}\)，对于给定已知的 \(y_{<t}\) 分布，由自回归特性，\(y_t\) 的分布也会随之确定，\(Y=\{y_1,\ldots,y_t\}\) 的分布也就已知。因此可以改写为：

\[
D_{\mathrm{KL}}(P(Y)\Vert Q(Y))
=
\sum_{t=1}^{n}
\mathbb{E}_{y_{<t}\sim P(y_{<t})}
\left[
D_{\mathrm{KL}}\left(P(y_t\mid y_{<t})\Vert Q(y_t\mid y_{<t})\right)
\right].
\]

这意味着，两个序列分布的KL散度，等于每一步条件概率分布之间KL散度的期望之和（期望是基于分布 \(P\) 生成的历史轨迹计算的）。

\FloatBarrier
\Needspace{5\baselineskip}
\hypertarget{ux53c2ux8003ux6587ux732e-107}{%
\subsection{参考文献}\label{ux53c2ux8003ux6587ux732e-107}}

\vspace{0.25em}
\begin{itemize}
\tightlist
\item
  Kullback, S., \& Leibler, R. A. (1951). \href{https://doi.org/10.1214/aoms/1177729694}{On Information and Sufficiency}. The Annals of Mathematical Statistics, 22(1), 79-86.
\item
  Lin, J. (1991). \href{https://doi.org/10.1109/18.61115}{Divergence measures based on the Shannon entropy}. IEEE Transactions on Information Theory, 37(1), 145-151.
\end{itemize}

\clearpage
